\color{unimain}\section[Основные функции]{Основные функции}\label{sec:funcs}
\color{black}

\needspace{4\baselineskip}
\color{unimain}{\subsection{Общие сведения}}
\color{black}

\begin{enumerate}[label=\arabic{section}.\arabic{subsection}.\arabic*, labelsep=4pt, leftmargin=0pt, itemindent=57pt]
    
\item
При вводе в работу устройство постоянно отслеживает состояние подключенных к нему аналоговых и дискретных сигналов. Из значений, полученных АЦП, при помощи расчетов и, при необходимости, цифровой фильтрации выделяются среднеквадратичные и/или мгновенные величины (фазных токов, гармонических составляющих токов, фазных и междуфазных напряжений), необходимые для работы функций в составе устройства.
\item
Оперативный вывод всех функций в составе устройства из работы происходит при появлении активного сигнала на входе <<Вывод терминала>>.
\item
Значения задаваемых параметров для токов функций в таблицах параметров настройки приведены по отношению к номинальному значению тока аналогового входа и указываются в относительных единицах (о.е).
\item
Коэффициенты возврата максимальных измерительных органов заданы равными 0,95, коэффициенты возврата минимальных измерительных органов заданы равными 1,05 (если в описании функции не указано другое значение).
\item
Основные характеристики функциональных блоков:
\begin{itemize}
    \item погрешность измерительных органов тока и напряжения (в том числе ИО по симметричным составляющим) не более 2,5\%, на границе частот 45 Гц и 55 Гц не более 5\%;
    \item время срабатывания ИО тока при повышении тока скачком от 0 до 2Iуст. не более 40 мс;
    \item время возврата реле ИО тока при снижении тока скачком от 2Iуст. до 0 не более 40 мс;
    \item время срабатывания ИО напряжения при повышении напряжения скачком от 0 до 2Uуст не более 40 мс;
    \item время возврата ИО напряжения при сбросе напряжения скачком от 2Uуст до 0 должно быть не более 40 мс; 
    \item погрешность выдержки времени -- не более 40 мс при выдержках менее 1,5 с и не более 2\% от уставки при выдержках времени выше 1,5 с;
    \item погрешность выдержки времени зависимой времятоковой характеристики не более 40 мс при уставке меньше 0,6 с, не более 12,5\% при 1,1≤I/Iср<5, не более 7,5\% при 5≤I/Iср<10 и не более 5\% при I/Iср≥10.
\end{itemize}
Специфические требования к характеристикам функциональных блоков приведены в описаниях соответствующих функций в данном разделе.
\item
Функции устройства могут использовать ток НП, полученный путем непосредственного измерения или вычисленный как сумма токов трех фаз. Настройка может быть выполнена в разделе <<Конфигурация>> ПО <<ЮНИТ-Сервис>>.
\item
Общая схема взаимосвязей между описанными в данном разделе функциями, входными и выходными сигналами устройства приведена в приложении \ref{app:fsu}.
Сводный перечень сигналов, формируемых устройством в процессе работы, приведен в таблице \ref{appA:tbl1} приложения \ref{app:signals}.

\end{enumerate}